%% ma: L10-B13-0100
%% lop: 10
%% chuong: 5
%% bai: 13
%% dang: 10.13.13
%% loai: DS
%% muc: [NB, TH, TH, VD]
%% dap_an: "ĐĐSS"
%% nguon: "chuyen-de-cac-so-dac-trung-cua-mau-so-lieu-khong-ghep-nhom-toan-10-chuong-trinh-moi.pdf (GV. Phan Nhật Linh) (P03), Bài 02 Dạng 2 Phần II Câu 4"
%% kien_thuc: "Lớp 10 Bài 13 (số trung bình, trung vị, tứ phân vị, mốt)"
%% trang_thai: chinh-thuc
%% ghi_chu: "Tự duyệt theo yêu cầu giáo viên 22/10/2026."
\begin{ex}
Cho số liệu thống kê về sản lượng chè thu được trong $1$ năm (kg/sào) của $20$ hộ gia đình như sau:
\begin{center}
\begin{tabular}{|c|c|c|c|c|c|c|c|}\hline
Giá trị $x$ & $111$ & $112$ & $113$ & $114$ & $115$ & $116$ & $117$\\\hline
Tần số & $1$ & $3$ & $4$ & $2$ & $5$ & $4$ & $1$\\\hline
\end{tabular}
\end{center}
Xét tính đúng sai của các khẳng định sau:
\choiceTF{\True Mốt của mẫu số liệu là $M_o=115$}{\True Số trung bình của mẫu số liệu là $\overline x=114{,}15$}{Trung vị của mẫu số liệu là $M_e=114$}{Các tứ phân vị của mẫu số liệu là $Q_1=112{,}5$; $Q_2=114{,}5$; $Q_3=115{,}5$}
\loigiai{a) Đúng: $115$ có tần số lớn nhất ($5$).

b) Đúng: $\overline x=\dfrac{2283}{20}=114{,}15$.

c) Sai: tần số tích lũy $1$; $4$; $8$; $10$; $15$ nên $x_{10}=114$, $x_{11}=115$ và $M_e=114{,}5$.

d) Sai: $Q_1=\dfrac{x_5+x_6}2=113$ (không phải $112{,}5$); $Q_3=\dfrac{x_{15}+x_{16}}2=\dfrac{115+116}2=115{,}5$.}
\end{ex}
